% Independent mathematical derivation of incoming 02_RESIDUE_BOUNDARY_VALUE.md.
% Inclusion fragment. The fixed-minor receiving map is proved for every chart.
\begingroup
\section{An explicit reflected trace and both original boundary metrics}
\label{sec:RBV}

Keep the original fixed hypothetical simple quartet, its full arithmetic
unit, actual period, logarithm branch and conductor. Write
\[
\begin{gathered}
a\equiv1\pmod4,\quad a\ge17,\quad c=a/2,\quad
q=(a+1)^2,\quad q'=(a-7)^2,\quad m=16a-48-v,\quad0\le v\le80,\\
\lambda_{u,b}=c+(2u-a)\delta+i(2b-a)\gamma,\qquad
0<\delta<\tfrac12,\quad\gamma>2,\quad0\le u,b\le a,\\
\chi(S)=\prod_{u,b=0}^a(S-\lambda_{u,b}),\quad
E=\mathbb C[S]/(\chi),\quad
\mathcal P_{q-1}\xrightarrow{\ \sim\ }E .
\end{gathered}\tag{RBV1}
\]
The last map sends an actual polynomial to its residue class, with
the unique degree-below-$q$ remainder as inverse. Roots are distinct
by equality of real and imaginary parts. No packet factors or coordinate
phases are discarded.

\subsection{The explicit polynomial and its exact trace}
Put
\[
t_0=\frac\pi{2\gamma},\qquad r_0=t_0\delta,\qquad
c_0=8\log4-10>0,\qquad
P_a(S)=\sum_{j=0}^{4a}\frac{t_0^{2j}(S-c)^{2j}}{(2j)!}.
\tag{RBV2}
\]
Its degree is $8a\le m-1$ since $8a\ge129\ge49+v$.
Also $8a<q$. Define the original finite reflected form
\[
\mathcal W(f,g)=\sum_{\lambda}\overline{f(a-\overline\lambda)}g(\lambda).
\tag{RBV3}
\]
The reflection permutes the complete grid and is an involution.
Changing its index in the conjugate of $\mathcal W(g,f)$ proves
that this form is Hermitian. The polynomial $P_a$ has real even
coefficients in $S-c$, so
$\overline{P_a(a-\overline\lambda)}=P_a(\lambda)$.
Hence $\mathcal W(P_a,P_a)=\sum_\lambda P_a(\lambda)^2$, a real
number because conjugation also permutes the complete grid.

For $z=\lambda_{u,b}-c$, the actual comparison value is
\[
\cosh(t_0z)=i\sin\!\left(\frac{\pi(2b-a)}2\right)
                   \sinh(r_0(2u-a)).\tag{RBV4}
\]
The sine is $+1$ or $-1$, and $2u-a$ is a nonzero odd integer.
Moreover
$t_0|z|\le at_0\sqrt{\delta^2+\gamma^2}<2a$:
$\delta/\gamma<1/4$ and $\pi<22/7$ give the final strict bound.
For $x\ge0$ and an integer $n\ge1$, the exponential series gives
$\sum_{j\ge n}x^j/j!\le x^ne^x/n!$, because
$(n+j)!\ge n!j!$. The elementary integral comparison
$\log n!=\sum_{j=1}^n\log j\ge\int_1^n\log x\,dx
=n\log n-n+1$ gives $n!\ge(n/e)^n$. Therefore
\[
\begin{split}
|P_a(c+z)-\cosh(t_0z)|
&\le\frac{e^{2a}(2a)^{8a+1}}{(8a+1)!}\\
&\le e^{2a}(e/4)^{8a+1}\le e^{-c_0a}=:\epsilon_a.
\end{split}\tag{RBV5}
\]
This bounds the full omitted even tail by a larger exponential tail.
Impose the explicit eventual guard $\epsilon_a\le r_0/4$.
Since $|\sinh(r_0(2u-a))|\ge r_0$, write the actual value as
$is+e$ with $s$ real and $|e|\le|s|/4$. Direct expansion gives
\[
\Re(is+e)^2\le-s^2+2|s||e|+|e|^2\le-7s^2/16.
\tag{RBV6}
\]
The elementary finite geometric sum gives exactly
\[
\begin{gathered}
\sum_{u=0}^a\cosh(2r_0(2u-a))
=\frac{\sinh(2r_0(a+1))}{\sinh(2r_0)},\\
\sum_{\lambda}\cosh(t_0(\lambda-c))^2
=-\frac{a+1}{2}
 \left[\frac{\sinh(2r_0(a+1))}{\sinh(2r_0)}-(a+1)\right].
\end{gathered}\tag{RBV7}
\]
Consequently the incoming finite sign and its constants are correct:
\[
\boxed{\mathcal W(P_a,P_a)\le-N_a<0,\qquad
N_a=\frac{7(a+1)}{32}
\left[\frac{\sinh(2r_0(a+1))}{\sinh(2r_0)}-(a+1)\right].}
\tag{RBV8}
\]
The bracket is $2\sum_u\sinh^2(r_0(2u-a))>0$.
There is a stronger explicit finite lower bound. For real $s$ and
$0\le\epsilon\le|s|/2$, maximizing $\Re(is+x+iy)^2$ over
$x^2+y^2\le\epsilon^2$ gives $-(|s|-\epsilon)^2$.
For $s>0$, substitute $x^2\le\epsilon^2-y^2$: the remaining
concave quadratic $-s^2-2sy+\epsilon^2-2y^2$ takes its maximum
on $[-\epsilon,\epsilon]$ at $y=-\epsilon$ because its vertex
is $-s/2\le-\epsilon$. The case $s<0$ follows by reflection.
Consequently RBV8 also holds with
\[
\begin{split}
N_a^{\rm sharp}
&=(a+1)\sum_{u=0}^a
 (|\sinh(r_0(2u-a))|-\epsilon_a)^2\\
&=(a+1)\left\{\frac12
 \left[\frac{\sinh(2r_0(a+1))}{\sinh(2r_0)}-(a+1)\right]
-2\epsilon_a\frac{\cosh(r_0(a+1))-1}{\sinh r_0}
 +(a+1)\epsilon_a^2\right\}.
\end{split}\tag{RBV8a}
\]
The additional finite sum follows by summing the positive and
negative odd-index hyperbolic sines separately. In particular the
coefficient $7/32$ may be strengthened to $9/32$ on the same guard.
The original smaller $N_a$ is retained in the downstream displayed
bounds, all of which also hold with $N_a^{\rm sharp}$.
The exact error from RBV7 is at most
\[
2\epsilon_a(a+1)\sum_{u=0}^a|\sinh(r_0(2u-a))|
                         +q\epsilon_a^2.\tag{RBV9}
\]
Since the sum is bounded by $(a+1)e^{r_0a}$, this error divided by
the magnitude in RBV7 is
$O_{\delta,\gamma}(a e^{-(c_0+r_0)a})$.
Keeping also the lower exponential and $(a+1)$ term in RBV7 proves
\[
\mathcal W(P_a,P_a)
=-\frac{(a+1)e^{2r_0(a+1)}}{4\sinh(2r_0)}[1+o(1)].
\tag{RBV10}
\]

\subsection{Original arithmetic density and finite full-form constants}
The source remains
\[
\begin{gathered}
v_h(s)=2\xi(s)/h(s),\quad
w_h(y)=|v_h(1/2+iy)|^2/(2\pi),\quad
d\mu_a(y)=w_h^{*a}(y)\,dy,\\
M_h(\theta)=\int_{\mathbb R}e^{\theta y}w_h(y)\,dy,
\quad\alpha=\pi/2,\quad
d\sigma(y)=|\Gamma(1/4+iy/2)|^2dy/(2\pi),\quad
\sigma(\mathbb R)=\sqrt{2\pi}.
\end{gathered}\tag{RBV11}
\]
RTM1--3/GTM1 proves evenness, positivity almost everywhere and the
finite transforms at $|\theta|\le\alpha$, using the complete unit.
For the simple quartet, the original global constants in GTM1 are
\[
\begin{gathered}
B=50,\quad p=7/2,\quad
\vartheta_h=\int_{-1}^1w_h(y)\,dy,\quad\mathfrak M=M_h(\alpha),\\
b_a=\frac{c_h\vartheta_h^{a-3}e^{-\alpha(a-3)}(a-2)^{-50}}
 {C_\Gamma2^{25}},\qquad
A_a=\frac{C_h^{\rm tilt}}{c_\Gamma}a^{p+1}\mathfrak M^{a-1},\\
b_a(1+y^2)^{-25}d\sigma\le d\mu_a\le A_a d\sigma,\qquad
c_\Gamma=\sqrt2e^{-7/6},\quad C_\Gamma=\sqrt{2\sqrt5}e^{2/3}.
\end{gathered}\tag{RBV12}
\]
The constants $c_h,C_h^{\rm tilt}$ are exactly the original
TW7/BT bounds retained in RTM2, including both half-lines. This proof
does not posit a positive lower bound for the one-factor zeta numerator.

Here is the finite polynomial comparison needed for the whole
boundary bands. In the original orthonormal Gamma frame,
multiplication by $y$ has adjacent entries
$\sqrt{(n+1)(n+1/2)}\le n+1$ (GTM7).
On polynomials of degree at most $D$, its two shifts each have norm
at most $D+1$. Thus
$\|yF\|_\sigma\le2(D+1)\|F\|_\sigma$.
For $F\ne0$, divide the positive integral $|F|^2d\sigma$ by its
own total only to apply the following scalar inequality; the source
measure itself is never changed. Convexity of $(1+x)^{-25}$ on
$x\ge0$, proved by its positive second derivative, gives its tangent
lower bound at the mean of $y^2$. Integration therefore gives
\[
\|F\|_{\mu_a}^2\ge
b_a[1+4(D+1)^2]^{-25}\|F\|_\sigma^2.
\tag{RBV13}
\]
For all degrees through $2q$ we may use the exact constants
\[
L_a^-=b_a[1+4(2q+1)^2]^{-25},\qquad L_a^+=A_a,
\qquad\log(L_a^+/L_a^-)=O_h(a+\log q).\tag{RBV14}
\]
The original Gamma recurrence and mass dictionary are the
Meixner--Pollaczek formulas in DLMF (18.19.8)--(18.19.9), (18.22.8)
\cite{human:dlmf18}, with the exact substitution and mass verified
in LT2 and NTG3. This is a full coefficient-form comparison, so it
applies before either original boundary minimization.

\subsection{Concentration of the complete relation bands}
Set $Y=2^{-16}q$ and $R=a\sqrt{\delta^2+\gamma^2}$, and retain the
explicit eventual guard $R\le Y/4$.
For every polynomial $P$ of degree at most $q$ and
$F(S)=\chi(S)P(S)$, the original roots give
\[
|\chi(c+iy)|\le(Y+R)^q\quad(|y|\le Y),\qquad
|\chi(c+iy)|\ge(q-R)^q\quad(q\le y\le2q).
\tag{RBV15}
\]
Use the shifted Legendre basis
$\phi_d(y)=\sqrt{(2d+1)/q}\,L_d(2y/q-3)$ on $[q,2q]$,
$0\le d\le q$. Its Rodrigues formula and squared norm are the
Legendre entries of DLMF (18.5.5), Tables 18.5.1 and 18.3.1
\cite{human:dlmf18}. The full finite identity
\[
L_d(z)=2^{-d}\sum_{j=0}^d\binom dj^2(z-1)^{d-j}(z+1)^j
\]
gives $|L_d(z)|\le10^d$ on $|z|\le4$, using
$\sum_j\binom dj^2=\binom{2d}{d}\le4^d$.
For $|y|\le Y$, expansion of the original polynomial
$P(c+iy)$ in this basis and Cauchy--Schwarz prove
\[
|P(c+iy)|^2\le\frac{(q+1)^2}{q}100^q
                         \int_q^{2q}|P(c+it)|^2dt.
\tag{RBV16}
\]
Indeed the basis-square sum is bounded by
$q^{-1}100^q\sum_{d=0}^q(2d+1)=q^{-1}100^q(q+1)^2$.
The lower Gamma density on $[q,2q]$ is
$c_\Gamma e^{-\pi q}(1+2q)^{-1/2}$. Multiplying RBV15--16,
integrating the full central interval using mass $\sqrt{2\pi}$,
and using RBV14 in the denominator proves
\[
\begin{gathered}
\int_{|y|\le Y}|F(c+iy)|^2d\mu_a
 \le\bar\eta_a\|F\|_{\mu_a}^2,\qquad
\bar\eta_a=\min(1,\eta_a),\\
\eta_a=\frac{L_a^+}{L_a^-}\frac{\sqrt{2\pi}}{c_\Gamma}
 \frac{(q+1)^2\sqrt{1+2q}}q e^{\pi q}100^q
 \left(\frac{Y+R}{q-R}\right)^{2q},\\
\log\eta_a\le-c_1q+O_h(a+\log q),\qquad
c_1=30\log2-\pi-\log100>0.
\end{gathered}\tag{RBV17}
\]
The last ratio is at most $2^{-15}$ on the stated guard. This
bound applies to every vector of either complete relation band,
including arbitrary linear combinations and all original cross terms.

\subsection{Both boundary energies, including strict positivity}
Let $e_j=\chi\pi_j/\sqrt{\nu_j}$, where $\pi_j$ is the actual
monic polynomial orthogonal for the full weight
$|\chi(c+iy)|^2d\mu_a$ in the original coordinate. Define
\[
\mathcal B_0=\operatorname{span}(e_0,\ldots,e_{q-1}),\qquad
\mathcal B_1=\operatorname{span}(e_1,\ldots,e_q),\qquad
\Delta_iP=-\operatorname{Proj}_{\mathcal B_i}P.
\tag{RBV18}
\]
These are the original COB two windows. Each band consists of
$\chi$ times polynomials of degree at most $q$, so RBV17 applies
to every unit vector in it.
Their exact link to the original minimum forms follows without a
choice of diagonal trial metric. Let $s_N:E\to\mathcal P_N$ send
a class to its unique minimum-norm representative in $d\mu_a$.
For the original representative $P\in\mathcal P_{q-1}$, its entire
affine fibre is $P+\chi\mathcal P_{N-q}$. Completing its orthogonal
square gives
$s_NP=P-\operatorname{Proj}_{\chi\mathcal P_{N-q}}P$.
Consequently
$\Delta_0=s_{2q-1}-s_{q-1}$ and
$\Delta_1=s_{2q}-s_q$, exactly as in RBV18.
Writing $G_N=s_N^*s_N$ in the unchanged class frame, orthogonality
of these two nested projection differences gives
$G_{q-1}-G_{2q-1}=\Delta_0^*\Delta_0$ and
$G_q-G_{2q}=\Delta_1^*\Delta_1$.
Thus each coefficient matrix is exactly
$\Gamma_{i,I}=L_I^*\Delta_i^*\Delta_iL_I$ on its stated original
residue injection. This proves the map into both boundary energies,
including all cross terms of the actual minimum.
These are precisely COB.5--6 and COB.13 in
\path{COUPLED_ORIGINAL_CONE_BOUNDARY.tex}. The general positive
coefficient-square operation is the Schur minimum in Boyd and
Vandenberghe \cite[Appendix A.5.5, p.~650]{human:boyd-vandenberghe};
the preceding formulas prove its full original-fibre specialization.

In fact both $\Delta_i$ are injective on the entire original
$\mathcal P_{q-1}$. Pair the roots with indices $u$ and $a-u$.
Since $(a+1)^2/2$ is even, their original factors give exactly
\[
\chi(c+iy)=\prod_{b=0}^a\prod_{0\le u<a/2}
\left[(y-(2b-a)\gamma)^2+(2u-a)^2\delta^2\right]>0.
\tag{RBV19}
\]
If $P\perp\mathcal B_0$, test its orthogonality against
$\chi P\in\mathcal B_0$ to obtain
$\int\chi(c+iy)|P(c+iy)|^2d\mu_a=0$, hence $P=0$.
If $P\perp\mathcal B_1$, choose the unique $b$ such that
$F=P-b\chi$ is orthogonal to $e_0$. Since $\chi$ is a multiple
of $e_0$, $F$ remains orthogonal to $e_1,\ldots,e_q$.
Its degree is at most $q$, so $\chi F$ belongs to
$\operatorname{span}(e_0,\ldots,e_q)$. Testing against it and using
RBV19 gives $F=0$. Then $P=b\chi$ and $\deg P<q$ force $P=b=0$.
This proves strict positivity without an assumed missing metric bound.

On the original source line,
$|P_a(c+iy)|\le\cosh(t_0|y|)\le e^{t_0|y|}$.
The exact convolution transform, with all masses retained, is
\[
\int e^{\theta y}d\mu_a=M_h(\theta)^a.
\]
Evenness and $2t_0<\alpha$ therefore give, with
$\theta_1=(\alpha-2t_0)/2>0$, $\theta_2=2t_0+\theta_1<\alpha$,
\[
\|P_a\|_{\mu_a}^2\le2M_h(2t_0)^a,\qquad
\|1_{|y|>Y}P_a\|_{\mu_a}^2
 \le2e^{-\theta_1Y}M_h(\theta_2)^a.
\tag{RBV20}
\]
For a unit $f\in\mathcal B_i$, split its inner product with $P_a$
over the two original integration regions. Cauchy--Schwarz and RBV17
give
$|\langle f,P_a\rangle|\le\sqrt{\bar\eta_a}\|P_a\|+
\|1_{|y|>Y}P_a\|$. Taking the supremum over the full unit sphere
of the band, then $(x+y)^2\le2x^2+2y^2$, proves
\[
\boxed{0<\|\Delta_iP_a\|_{\mu_a}^2\le B_a,
\quad B_a=4\bar\eta_aM_h(2t_0)^a+
4e^{-\theta_1 2^{-16}q}M_h(\theta_2)^a,\quad i=0,1.}
\tag{RBV21}
\]
No band-dimension multiplier enters this estimate.

\subsection{Relative eigenvalue and evaluated discrepancy}
Let $G_i^E$ be the Gram of $\Delta_i$ on the actual monomial frame
of $\mathcal P_{q-1}$; it is strictly positive by RBV19.
Let $W^E$ be the matrix of RBV3 in that same frame. Then
\[
T_i^E=(G_i^E)^{-1/2}W^E(G_i^E)^{-1/2}
\]
is Hermitian. The coefficient column of $P_a$ gives the actual
Rayleigh quotient, so
\[
\lambda_{\min}(T_i^E)\le-N_a/B_a<0,\qquad
\|I-T_i^E\|\ge1+N_a/B_a.\tag{RBV22}
\]
This follows by diagonalizing the Hermitian matrix and expressing
the Rayleigh quotient as a positive weighted average of its real
eigenvalues. With the original two matrices, the discrepancy is
\[
\begin{split}
\|\Delta_iP_a\|_{\mu_a}^2-\mathcal W(P_a,P_a)
&=\frac{(a+1)e^{2r_0(a+1)}}{4\sinh(2r_0)}[1+o_h(1)],\\
\lim_{a\to\infty}\frac{\|\Delta_iP_a\|^2-\mathcal W(P_a,P_a)}
 {a e^{\pi\delta a/\gamma}}
&=\frac{e^{\pi\delta/\gamma}}{4\sinh(\pi\delta/\gamma)}.
\end{split}\tag{RBV23}
\]
Indeed RBV21 has logarithm at most
$-c_*q+O_h(a+\log q)$, where
\[
c_*=\min\{30\log2-\pi-\log100,
                   2^{-16}(\pi/2-\pi/\gamma)/2\}>0.
\tag{RBV24}
\]
This is negligible compared with RBV10. Also
$\log N_a=2r_0a+\log a+O_h(1)$, so RBV22 proves
\[
\liminf_{a\to\infty}q^{-1}\log[-\lambda_{\min}(T_i^E)]\ge c_*.
\tag{RBV25}
\]
All limits run over the original $a\equiv1\pmod4$ family at fixed
actual quartet/unit and eventually satisfy the explicitly stated
guards. No actual-period numerical eigenvalue is invented.

\subsection{Exact return through every original conductor chart}
This paragraph records the precise map needed to receive these
results on the supplied COB coefficient image, without identifying
two different sections of the same quotient. Write
$\chi(S)=\sum_{j=0}^qc_jS^j$, $c_q=1$, and define
\[
\mathfrak p_n(S)=\sum_{j=n+1}^qc_jS^{j-n-1},\qquad
\mathbb V_{n,\lambda}=\lambda^n,\qquad
\mathcal L\nu=\sum_\lambda\nu_\lambda\frac{\chi(S)}{S-\lambda},
\quad F=\mathcal L\mathbb V^{-1}.
\tag{RBV26}
\]
NCB5--7 proves $Fe_n=\mathfrak p_n$ by the coefficient of $S^{-1}$
in $\mathfrak p_nS^j/\chi$. Thus $F$ is invertible and its inverse is
the exact Laurent map
$(F^{-1}P)_n=[S^{-n-1}]P(S)/\chi(S)$, $0\le n<q$.
Let $C^{\mathsf T}$ be the original conductor multiplication in
exponential coefficients, and write
$\mathbb VC^{\mathsf T}=P_vJ$ with the original $v$ leading zeros.
For every retained invertible $q'$-row minor $J_I$, define
\[
\begin{gathered}
H_I=J_I^{-1}E_I^{\mathsf T},\quad
Q_I=E_{I^c}^{\mathsf T}-J_{I^c}H_I,\quad
L_I=FP_vE_{I^c},\quad C_E=\mathcal LC^{\mathsf T},\\
\eta_I=H_IP_v^{\mathsf T}F^{-1}P_a,\quad
z_I=Q_IP_v^{\mathsf T}F^{-1}P_a,\qquad
\boxed{P_a=C_E\eta_I+L_Iz_I.}
\end{gathered}\tag{RBV27}
\]
There is no low component because $\deg P_a<m$ makes all Laurent
coordinates below $q-m\ge v$ zero. To prove the displayed identity,
check the $I$ and $I^c$ rows to get
$JH_I+E_{I^c}Q_I=I$. Multiply by $FP_v$ and use
$FP_vJ=C_E$. These are the exact original residue and conductor maps;
no coefficient section is inserted in place of another.

In particular the consecutive minor
$I_0=\{0,\ldots,q'-1\}$ has determinant
\[
\det J_{I_0}=\mu_v^{q'}
 \prod_{r=0}^{q'-1}\binom{v+r}{v}\det\mathbb V_{\rm lower}\ne0.
\tag{RBV28}
\]
Indeed its row $r$ is the lower-root evaluation of a degree-$r$
polynomial with leading coefficient $\binom{v+r}{v}\mu_v$;
triangular coefficient elimination gives exactly this determinant.
Here $\eta_{I_0}=0$ and
$\operatorname{im}L_{I_0}=\mathcal P_{m-1}$ by the distinct monic
degrees $m-1,m-2,\ldots,0$ of its columns. The literal inverse is
\[
(z_{I_0})_j=[S^{-(q-m+j+1)}]P_a(S)/\chi(S).
\tag{RBV29}
\]
Therefore RBV8--25 apply to the Gram
$\Gamma_{i,I_0}=L_{I_0}^*G_i^EL_{I_0}$ and reflected matrix
$L_{I_0}^*W^EL_{I_0}$ with precisely the same bounds, because the
Rayleigh vector is the explicit column RBV29. For any actual other
chart $I$, RBV27 supplies the full original coefficient vector
$(\eta_I,z_I)$ instead. With $M_I=[C_E,L_I]$, the same calculation
holds on its full image with the two exact matrices
\[
M_I^*G_i^EM_I=\begin{pmatrix}
C_E^*G_i^EC_E&C_E^*G_i^EL_I\\
L_I^*G_i^EC_E&L_I^*G_i^EL_I
\end{pmatrix},\qquad M_I^*W^EM_I.
\tag{RBV30}
\]
The columns of $M_I$ form the original direct conductor/native
decomposition in the high-moment space, hence this Gram is positive
definite. All cross terms in RBV30 are essential.
This gives a proved negative class and exponential relative bound
through every complete original chart. Selecting only the native
block in a nonconsecutive chart is a further, different restriction:
RBV27 evaluates its omitted conductor component exactly instead of
declaring it zero. The companion original-image derivation identifies
which selector is the specified current COB map; no root ordering
alone is used here as a proof of an ordering of row subsets.
For the original four-copy maps use $I_4\otimes L_{I_0}$ or
$I_4\otimes M_I$ and the coefficient column with the displayed
one-copy preimage in its first block and zero in the other three.
Its reflected value and each boundary energy are exactly the
one-copy values. Hence RBV22--25 have the same constants on these
original four-copy images, with no additional rank factor.

\subsection{The complete tensor residue, its multiplier, and the residual}
We now expand the full FTR/FRT comparison on the present original
simple-quartet domain. For the four original roots
$z_{\epsilon\eta}=1/2+(2\epsilon-1)\delta+i(2\eta-1)\gamma$,
retain every original unit value in
\[
w_{\epsilon\eta}=\frac{v_h(z_{\epsilon\eta})}{h'(z_{\epsilon\eta})}\ne0,
\quad p(X,Y)=\sum_{\epsilon,\eta=0}^1w_{\epsilon\eta}X^\epsilon Y^\eta,
\quad c_{u,b}^{\rm ten}=[X^uY^b]p(X,Y)^a.
\tag{RBV31}
\]
There are no omitted jets: the one-factor multiplicity is exactly
one here. The original full-unit residue formula FTR.8, by taking
the simple residues in every ordered tensor factor, gives
\[
\mathcal B(f,g)=\sum_{u,b=0}^a
c_{u,b}^{\rm ten} f^{\sharp_a}(\lambda_{u,b})g(\lambda_{u,b}).
\tag{RBV32}
\]
Indeed each ordered root tuple contributes the product of its
$a$ unchanged $w_{\epsilon\eta}$ factors; the exponents of $X,Y$
record the two original plus counts. Expansion of $p^a$ counts
every ordered tuple once. Thus the tensor pullback multiplier in
the cyclic residue form is the actual class with values
$b_{\rm ten}(\lambda_{u,b})=\chi'(\lambda_{u,b})c_{u,b}^{\rm ten}$.
This is the complete simple-root specialization of FTR.6--12.

Put $\mathcal D=\{(u,b):c_{u,b}^{\rm ten}\ne0\}$ and
$\mathcal Z=\{0,\ldots,a\}^2\setminus\mathcal D$, and retain the
original idempotents
$e_{u,b}=\chi(S)/((S-\lambda_{u,b})\chi'(\lambda_{u,b}))$.
The trace multiplier and complementary form are exactly
\[
\begin{gathered}
a_{\rm tr}=\sum_{(u,b)\in\mathcal D}
              (c_{u,b}^{\rm ten})^{-1}e_{u,b},\qquad
\mathcal W_{\mathcal Z}(f,g)=\sum_{(u,b)\in\mathcal Z}
              f^{\sharp_a}(\lambda_{u,b})g(\lambda_{u,b}),\\
b_{\rm ten}a_{\rm tr}=e_{\mathcal D}\chi'\quad\text{in }E,
\qquad
\boxed{\mathcal W(f,g)=\mathcal B(f,a_{\rm tr}g)
                              +\mathcal W_{\mathcal Z}(f,g).}
\end{gathered}\tag{RBV33}
\]
All equations follow by evaluation at every simple root; the
evaluation map is an isomorphism, so they prove the typed algebra
map, not just matching dimensions. They are exactly FRT.2--5.
The trace equality can also be verified by taking the residue of
$\chi'f^{\sharp_a}g/\chi$ at each root. The classical multiplication
trace construction is recorded by the Stacks Project Authors,
Section 49.3, tags 0BVH and 0BJF \cite{human:stacks-trace}; the
displayed root computations prove the full specialization here.

In particular the retained residual has the exact finite value
\[
\mathcal W_{\mathcal Z}(P_a,P_a)
 =\sum_{(u,b)\in\mathcal Z}P_a(\lambda_{u,b})^2.
\tag{RBV34}
\]
Conjugation of the full unit and $h$ gives
$w_{\epsilon,1-\eta}=\overline{w_{\epsilon\eta}}$;
their original reflection gives
$w_{1-\epsilon,\eta}=-\overline{w_{\epsilon\eta}}$.
Since $a$ is odd, the coefficient sets $\mathcal D,\mathcal Z$
are closed under both grid symmetries. Thus both partial sums of
$P_a(\lambda)^2$ are real. RBV6 proves that RBV34 is nonpositive,
and strictly negative whenever $\mathcal Z$ is nonempty. This
does not determine that set by a rank assumption: its exact definition
retains all cancellations in the actual coefficients RBV31.

There is also a negative pairing directly in the retained tensor
residue without needing to determine its interior zero coefficients.
The four edge coefficients of $p^a$ are binomial products of
nonzero original $w$'s, so every perimeter point belongs to
$\mathcal D$. In particular its two complete edges $u=0,a$
contain $2(a+1)$ points with absolute comparison value
$\sinh(r_0a)$. Applying RBV8a at these points and RBV6 at all
remaining points of $\mathcal D$ proves
\[
\boxed{\mathcal B(P_a,a_{\rm tr}P_a)
 \le-2(a+1)(\sinh(r_0a)-\epsilon_a)^2<0.}\tag{RBV35}
\]
The second argument is the explicitly computed full polynomial
class $a_{\rm tr}P_a\bmod\chi$; it is not identified with a native
coefficient vector unless its Laurent coordinates satisfy that
specific section's membership test.

Finally the original arithmetic injection is
$\mathcal I:E\hookrightarrow Q^{\otimes a}$ with the complete
$[v_h]_h^{\otimes a}$ factor. Its inverse on its named image is
the original full Mellin-jet observation, division by that invertible
Taylor factor and inverse sum substitution. Thus
$\mathcal I M_{a_{\rm tr}}\mathcal I^{-1}$ is the exact arithmetic
map realizing RBV33--35 on $\mathcal I(e_{\mathcal D}E)$.
In RBV35 the first argument may be written $e_{\mathcal D}P_a$,
since the complementary first-input block pairs to zero under
$\mathcal B$ by RBV32; the separate value RBV34 remains in
the full reflected trace.
It commutes with the original sum action because both act by
multiplication before this injection. The proper boundary image
uses instead the full COB.8 map
$\Psi_2:[P]\mapsto[P(D_1+\cdots+D_a)F_h^{\otimes a}]$
on $\mathbb C[S]/\chi^{a+1}$, whose one-factor Taylor unit is
$[v_h]_{h^2}$; these two different original target spaces are not
silently identified. The ordinary class is recovered on the named
boundary image by $\mathcal I\Delta_i^{-1}\Psi_2^{-1}$ exactly
as in COB.12. This retains the full tensor-residue residual and
every original unit. The finite packet trace is not substituted
for the full infinite Weil distribution, and none of these values
is assigned to the native minimum metric.

The tensor calculation also has a relative-metric receiver, with
the residual retained at matrix level. On any of the certified
images above let $L$ denote its displayed injection and set
$\Gamma_i=L^*G_i^EL$. With $W_{\mathcal D}^E,W_{\mathcal Z}^E$
the two explicitly computed reflected matrices, define
\[
\begin{gathered}
T_{i,\mathcal D}=\Gamma_i^{-1/2}L^*W_{\mathcal D}^EL\Gamma_i^{-1/2},
\qquad T_{i,\mathcal Z}=\Gamma_i^{-1/2}L^*W_{\mathcal Z}^EL\Gamma_i^{-1/2},\\
T_i=T_{i,\mathcal D}+T_{i,\mathcal Z},\qquad
\boxed{\lambda_{\min}(T_{i,\mathcal D})
\le-\frac{2(a+1)(\sinh(r_0a)-\epsilon_a)^2}{B_a}<0.}
\end{gathered}\tag{RBV36}
\]
Both partial matrices are Hermitian by their proved symmetry sets.
Use the very same preimage of $P_a$ as in RBV22 and the strict
boundary bound RBV21 to prove the displayed inequality. Its
negative magnitude again has logarithmic lower rate $c_*$ after
division by $q$. This is a bound for the actual tensor-residue
comparison $\mathcal B(f,a_{\rm tr}g)$; the other matrix remains
in the exact displayed sum. No equality of the smallest eigenvalues
of those three matrices is asserted.
\endgroup


% Complete proof fragment. Bibliography entries are in REFERENCES.tex.
\begingroup
\section{Uniform recovery on the original fixed-period observation image}
\label{sec:UPR}

This calculation derives U1--U8 of the supplied
\emph{A uniform inverse bound for the original period observation},
19 September 2026.  It corrects one finite-module degree statement and
makes the linear division operator and its degree-uniform bound explicit.
The actual observation, Fourier frame and inherited Gram are AKS7--10
in the existing cumulative source \cite{programme:aks-current}.
Every period constant below is at the same fixed original nonzero
admitted period and its original logarithm branch.  No numerical
certificate for that period is claimed.

\subsection{The unchanged coefficient-dual map and its exact kernel}
Retain a simple stipulated quartet
\[
\tfrac12\pm\delta\pm i\gamma,\qquad
0<\delta<\tfrac12,\quad\gamma>2,
\]
the original arithmetic unit \(U\), primary idempotent matrix
\(\mathcal E_{\rm prim}\), and period matrix \(\Pi\).  The precise
one-factor matrix and Fourier data are
\[
F_{jt}=\zeta^{jt}-1,\quad
(F^{-1})_{tj}=\zeta^{-tj}/5,\quad
\zeta=e^{2\pi i/5},\quad
H=F^{-1}\Pi U\mathcal E_{\rm prim},\quad
D=\operatorname{diag}(\zeta,\zeta^2,\zeta^3,\zeta^4).
\tag{UPR1}
\]
The original provider establishes invertibility of \(H\).  In
\(\mathcal A=\mathbb C[x_1,x_2,x_3,x_4]\), put
\[
Q(x)=Q_0(H^{-1}x),\qquad Q_0(t)=t_1t_4-t_2t_3.
\]
For degree \(k\ge0\) let
\(\mathcal B_k=\operatorname{span}\{Z^aW^b:0\le a,b\le k\}\).
The original coefficient-dual map is
\[
\Theta_k:\mathcal A_k\longrightarrow\mathcal B_k,\qquad
f(x)\longmapsto f\bigl(H(ZW,Z,W,1)^{\mathsf T}\bigr).
\tag{UPR2}
\]
The arithmetic receiver below uses the original subsequence
\(k\equiv1\pmod4,\ k\ge5\); these finite polynomial constructions
also define the intermediate graded maps at every degree.

Here is an exact proof of its kernel and an explicit full lift.
For each pair \(0\le a,b\le k\), put \(n=\max(0,a+b-k)\), and define
\[
\ell_{a,b}(t)=
t_1^nt_2^{a-n}t_3^{b-n}t_4^{k-a-b+n}.
\tag{UPR3}
\]
All exponents are nonnegative, their sum is \(k\), and substitution
\(t=(ZW,Z,W,1)\) gives \(Z^aW^b\).  Every homogeneous monomial in
the \(t_i\) can be reduced by \(t_1t_4=t_2t_3\) until either its
first or its fourth exponent is zero.  The two exponents
\(a=n_1+n_2\), \(b=n_1+n_3\), and its total degree \(k\) then
determine all four exponents uniquely, yielding (UPR3).
Thus the remaining monomials have linearly independent images, and
\[
\ker\Theta_k=Q\mathcal A_{k-2},\qquad
\Theta_k\mathcal A_k=\mathcal B_k,
\tag{UPR4}
\]
where a negative homogeneous degree means zero.  This argument
also proves the complete relation ideal; it is not a dimension-only
identification.

Define the linear lift on the displayed biform frame by
\(L_k(Z^aW^b)=\ell_{a,b}(H^{-1}x)\).
For the ordinary coefficient \(\ell^1\) norm put
\[
h_-=\max\left(1,\max_i\sum_j|(H^{-1})_{ij}|\right).
\]
The sum of absolute coefficients of a product is at most the
product of those sums.  Apply this to the \(k\) linear factors of
(UPR3), and then to a linear combination of biform monomials:
\[
\Theta_kL_k=I_{\mathcal B_k},\qquad
\|L_kf\|_1\le h_-^k\|f\|_1.
\tag{UPR5}
\]
The full arithmetic unit and period remain in \(H\) and \(h_-\).

\subsection{A fixed homogeneous module with the corrected degree bound}
Put
\[
\mathcal S=\mathbb C[X_1,X_2,X_3,X_4],\quad X_i=x_i^5,\qquad
\mathcal R=\{0,1,2,3,4\}^4.
\]
The exact decomposition
\[
\mathcal A=\bigoplus_{r\in\mathcal R}\mathcal S x^r
\tag{UPR6}
\]
follows by Euclidean division of each monomial exponent by five.
The invariant submodule for the original \(D\)-action is
\[
\mathcal A^D=\bigoplus_{r\in\mathcal R_0}\mathcal S x^r,\qquad
\mathcal R_0=\{r:\textstyle\sum_{i=1}^4 i r_i=0\bmod5\}.
\tag{UPR7}
\]
There are \(625\) residue vectors and \(125\) invariant residue
vectors: the fourth residue is uniquely determined by the first
three in the displayed congruence.

Use the original grading \(\deg X_i=5\) and \(\deg x^r=|r|\).
For the second source summand below the residue basis has the
additional degree-two shift:
\[
\mathcal M:
\left(\bigoplus_{r\in\mathcal R_0}\mathcal Sx^r\right)
\oplus
\left(\bigoplus_{r\in\mathcal R}\mathcal Sx^r[2]\right)
\longrightarrow
\bigoplus_{r\in\mathcal R}\mathcal Sx^r,\qquad
(f,g)\longmapsto f+Qg.
\tag{UPR8}
\]
It is an exact degree-preserving map, and
\[
\operatorname{im}\mathcal M=\mathcal A^D+Q\mathcal A.
\tag{UPR9}
\]

The entries of the multiplication-by-\(Q\) block have total
\(X\)-degree at most \emph{two}, rather than the at-most-one assertion
in the incoming U2.  Indeed multiplication of \(x^r\) by a quadratic
monomial can cause at most two carries in division of exponents
by five.  For distinct \(i,j\), the term \(x_ix_j\) and a residue
with \(r_i=r_j=4\) produce the entry \(X_iX_j\).  The actual
grading and the fact that this is one fixed finite polynomial
matrix are unchanged by this correction.

\subsection{Finite basis, homogeneous lifts and a specified linear division}
We use Buchberger's polynomial-basis method
\cite{human:buchberger2006}, with its finiteness supplied by
Dickson's lemma \cite{human:dickson1913}.  The needed module argument
and norm estimates are supplied here.

Enumerate the \(625\) target residues once in lexicographic order,
with integer position \(b_r\in\{0,\ldots,624\}\).  Order target
monomials \(X^\alpha x^r\) lexicographically by
\((\alpha_1,\alpha_2,\alpha_3,\alpha_4,b_r)\).
This is a well-order compatible with multiplication by every
\(X^\beta\).

Start with the nonzero columns of (UPR8), retaining their corresponding
source coordinate vectors.  Make each column monic by division by
its actual nonzero leading coefficient.  For pairs whose leading
target position agrees, form the usual least-common-monomial
cancellation, reduce it using the current monic list, and append
each nonzero remainder with its simultaneously updated source
lift.  Order the finite pair lists lexicographically to specify
the choices.  A reduction of a homogeneous vector by a homogeneous
column retains its full shifted degree, because the cancelled
monomials occupy the same target position.  Every resulting column
\(g_j\) and retained lift \(h_j\) are therefore homogeneous and
\[
\mathcal M h_j=g_j.
\tag{UPR10}
\]

For completeness, an infinite sequence in \(\mathbb N^d\) has an
earlier term componentwise at most a later one.  For \(d=1\), either
some value occurs infinitely often, or an infinite strictly
increasing subsequence can be selected.  In either case there is
an infinite nondecreasing subsequence.  Apply this to the first
coordinate, and apply the induction hypothesis to the other
coordinates of that subsequence.  This proves the assertion for
every finite \(d\), including \(d=4\).
An infinite sequence of appended leading monomials would contain
infinitely many in one target position and then an earlier one
dividing a later one.  That contradicts the complete remainder
reduction before appending.  Hence only finitely many columns are
appended, and the pair algorithm terminates.

The output is a module Gr\"obner basis.  To verify this criterion
directly, take a nonzero element of \(\operatorname{im}\mathcal M\)
and a representation by monomial multiples of the output columns
whose largest occurring leading monomial is least possible; among
these, minimize the number of terms at that largest monomial.
If it exceeds the leading monomial of the represented vector,
at least two terms at the same target position cancel.
Their least-common-monomial combination is a multiple of one of
the processed pairs.  That pair has a representation with all
occurring leading monomials strictly smaller than its cancelled
monomial, because its remainder is zero.  Substitute this
representation and combine the remaining terms at the largest
monomial.  This decreases their number, or the largest monomial
itself, a contradiction.  Thus some leading basis monomial divides
the leading monomial of the given element.  Division consequently
has zero remainder exactly on the image module.
A reduced basis may be obtained by discarding divisible leading
monomials and reducing the tails, updating their lifts throughout;
reducedness is not needed in the subsequent estimate.

We now specify positive integer scores that strictly decrease in
every rewrite.  Let \(D_*\) be the maximum of all absolute exponent
differences, coordinate by coordinate, between a leading term and
a tail term in the finite list \(g_j\); set \(D_*=0\) when all tails
are empty.  Set
\[
T_*=2(D_*+625)+2,\quad
w_i=T_*^{\,5-i}\quad(1\le i\le4),\quad
W_*=\max_i w_i,\quad B_*=624.
\tag{UPR11}
\]
Give \(X^\alpha x^r\) score
\(\operatorname{sc}(\alpha,r)=\sum_iw_i\alpha_i+b_r\).
If a leading and tail term first differ at variable index \(i\),
their score difference is at least
\[
T_*^{5-i}
-D_*\sum_{j>i}T_*^{5-j}-624>0.
\]
The last inequality follows from
\(D_*/(T_*-1)<1/2\) and \(T_*/2>624\).
If all four exponents agree, the leading position has strictly
larger \(b_r\).  Thus every such difference is an integer at least
one.  Multiplication by a monomial shifts both scores equally.
All scores are nonnegative.

Set the following fixed constants from the finite basis and its
actual lifts:
\[
C_*=\max\left(2,1+\max_j\|\operatorname{tail}(g_j)\|_1\right),
\qquad H_*=\max\left(1,\max_j\|h_j\|_1\right).
\tag{UPR12}
\]
Here and below coefficient norms in the free module use its
displayed residue basis; under (UPR6) these are precisely the
ordinary polynomial coefficient norms.

Define a linear division operator by a recursion on individual
monomials, rather than by choices depending on a whole input
polynomial.  If no leading \(g_j\) divides a target monomial \(t\),
set \(\mathcal R(t)=0\), \(\operatorname{NF}(t)=t\).
Otherwise use the first such \(g_j\), put \(m=t/\operatorname{LM}g_j\),
and define
\[
\begin{split}
\mathcal R(t)&=m h_j-\mathcal R(m\,\operatorname{tail}g_j),\\
\operatorname{NF}(t)&=-\operatorname{NF}
                         (m\,\operatorname{tail}g_j).
\end{split}
\tag{UPR13}
\]
Extend each right side linearly over its finitely many monomials.
All recursive arguments have lower score, so this is a finite
definition.  Finally extend to arbitrary polynomials by their
monomial expansion.  The resulting maps are linear, and induction
in the recursion proves
\[
f=\mathcal M\mathcal R(f)+\operatorname{NF}(f).
\tag{UPR14}
\]
The normal form has only irreducible monomials; if \(f\) is in
the image module, (UPR14) and the just-proved leading-monomial
criterion imply that it is zero.  No choice of a degree-dependent
minor is involved.

If \(t\) has score \(s\), another induction proves
\[
\|\mathcal R(t)\|_1\le H_*C_*^{s+1}.
\tag{UPR15}
\]
Indeed the first lift has norm at most \(H_*\) and the total
absolute coefficient mass of its tails is at most \(C_*-1\).
Their scores are at most \(s-1\), so the induction gives the bound
\[
H_*+(C_*-1)H_*C_*^s\le H_*C_*^{s+1}.
\]
At score zero a reducible term has no tail, and the same inequality
holds directly.  Monomial multiplication does not change coefficient
norms.  Homogeneity of (UPR10) shows that this recursion preserves
the original shifted degree throughout.

On degree \(k\), \(5|\alpha|+|r|=k\), so the score is at most
\(W_*k/5+B_*\).  By linearity, for every degree-\(k\) element of
the actual image module,
\[
\mathcal M\mathcal R(f)=f,\qquad
\|\mathcal R(f)\|_1
\le H_*C_*^{W_*k/5+B_*+1}\|f\|_1.
\tag{UPR16}
\]
These constants are defined by one finite calculation over the
fixed field of actual coefficients.  The proof establishes their
finiteness and the uniform bound.  It does not claim that this
basis has been numerically computed for an unspecified hypothetical
quartet or that exact zero tests on arbitrary supplied complex
numbers are available to a numerical implementation.

\subsection{A linear right inverse on the original invariant image}
Let \(f\in\Theta_k(\mathcal A_k^D)\).  Its fixed lift \(L_kf\)
differs from an invariant lift by a multiple of \(Q\), by (UPR4).
Thus \(L_kf\in\operatorname{im}\mathcal M\).
Write \(\mathcal R(L_kf)=(p,g)\) in (UPR8).  Then
\[
p\in\mathcal A_k^D,\qquad L_kf=p+Qg,\qquad \Theta_kp=f.
\]
The map \(f\mapsto p\) is linear by (UPR5) and (UPR13).
Since \(\dim\mathcal B_k=(k+1)^2\),
\[
\|p\|_2\le\|p\|_1
\le\mathscr R_k\|f\|_2,\qquad
\mathscr R_k=
H_*C_*^{W_*k/5+B_*+1}h_-^k(k+1).
\tag{UPR17}
\]
In particular \(\log\mathscr R_k=O_{H}(k+\log(k+1))\).
Here \(k\) is the original tensor degree; the coefficient module
and all its constants are fixed.

To compare with the actual matrix, expand
\(\prod_t(\sum_\alpha H_{t\alpha}t_\alpha)^{\nu_t}\)
for each invariant \(\nu\).
The coefficient with plus-real count \(a\), plus-imaginary count
\(b\), is
\[
B_{\nu;(a,b)}
=\sum_{\substack{\sum_\alpha n_{t\alpha}=\nu_t\\
                  \text{original plus counts }a,b}}
\left(\prod_t\frac{\nu_t!}{\prod_\alpha n_{t\alpha}!}\right)
\prod_{t,\alpha}H_{t\alpha}^{\,n_{t\alpha}},
\tag{UPR18}
\]
exactly AKS8.  Thus the matrix of the invariant restriction of
(UPR2) is \(B_k^{\mathsf T}\) in these unchanged coefficient frames.
For any map \(T\), its least-norm preimage has norm at most the
preimage supplied by a right inverse.  Equation (UPR17) therefore
bounds the Moore--Penrose inverse of \(B_k^{\mathsf T}\) by
\(\mathscr R_k\).  A singular-value factorization
\(B_k=U\Sigma V^*\) gives
\(B_k^{\mathsf T}=\bar V\Sigma(\bar U)^*\), with exactly the
same nonzero singular values.  Consequently
\[
\|(B_k)^+\|\le\mathscr R_k.
\tag{UPR19}
\]
Zero singular values are kept as the actual observation kernel.
If its whole image is zero, the pseudoinverse is zero and the
same assertion holds.

For comparison the forward norm is also degree-exponential.
With \(h_+=\max(1,\max_i\sum_j|H_{ij}|)\), expansion of the same
linear factors has coefficient \(\ell^1\) norm at most \(h_+^k\).
Substitution \(t=(ZW,Z,W,1)\) only combines coefficients and cannot
increase that norm.  It follows that
\[
\|B_k\|=\|B_k^{\mathsf T}\|
\le h_+^k\sqrt{\binom{k+3}{3}}.
\tag{UPR20}
\]
Hence the condition number on its nonzero coefficient singular
values is \(\exp(O_H(k+\log(k+1)))\).

\subsection{The exact original word Gram, with each multiplicity once}
Direct summation of fifth roots of unity gives
\[
(F^*F)_{tu}
=\sum_{j=1}^4(\zeta^{-jt}-1)(\zeta^{ju}-1)
=5(\delta_{tu}+1).
\tag{UPR21}
\]
Thus \(F^*F=5(I+\mathbf1\mathbf1^*)\), whose eigenvalues are
five with multiplicity three and twenty-five with multiplicity one.

For \(|\nu|=k\), let \(e_\nu\) be the sum of all distinct ordered
standard tensor words with multiplicity vector \(\nu\).
These are the unscaled symmetric-word vectors of the original
provider.  Different \(e_\nu\) are orthogonal and
\[
\|e_\nu\|^2=k!/\nu!,\qquad
1\le k!/\nu!\le4^k.
\tag{UPR22}
\]
The upper bound follows because the sum of all these multinomial
integers is \(4^k\).
If \(\mathcal F_0\) has columns \(F^{\otimes k}e_\nu\) for the
original invariant index set, then its inherited Gram entries are
\[
(\mathcal F_0^*\mathcal F_0)_{\nu\mu}
=\sum_{\substack{\sum_u n_{tu}=\nu_t\\
                  \sum_t n_{tu}=\mu_u}}
\frac{k!}{\prod_{t,u}n_{tu}!}
\prod_{t,u}[5(\delta_{tu}+1)]^{n_{tu}}.
\tag{UPR23}
\]
Indeed expansion of
\((\bar x^{\mathsf T}F^*Fy)^k\) gives the displayed coefficients;
equivalently expand both ordered tensor-word sums.
This is the exact AKS10 formula, including every row cross term.

The squared singular values of \(F^{\otimes k}\) lie between
\(5^k\) and \(25^k\), since tensor eigenvalues multiply.
For an arbitrary coefficient column \(a=(a_\nu)\), (UPR22) gives
\[
\sum_\nu|a_\nu|^2
\le\left\|\sum_\nu a_\nu e_\nu\right\|^2
\le4^k\sum_\nu|a_\nu|^2.
\]
Apply the tensor singular-value bounds to this same vector,
restricted to the actual invariant indices:
\[
\boxed{5^kI\le\mathcal G_0:=\mathcal F_0^*\mathcal F_0
\le100^kI.}
\tag{UPR24}
\]
The row multiplicities in (UPR18) and the word Gram (UPR23) have
different exact roles.  No additional multinomial is inserted
into either of them.

The original AKS9 map in root-value coordinates is
\[
\jmath\Lambda_k v=\mathcal F_0B_kv.
\tag{UPR25}
\]
The map \(\mathcal F_0\) is injective, with
\(\|\mathcal F_0^+\|\le5^{-k/2}\) by (UPR24).
For \(y\) in the actual image of (UPR25), define
\[
T_k y=B_k^+\mathcal F_0^+y.
\]
Then \(B_kT_ky=\mathcal F_0^+y\) and
\(\mathcal F_0B_kT_ky=y\), because the latter vector lies in
\(\operatorname{im}B_k\).  Furthermore
\[
\|T_ky\|_2\le5^{-k/2}\mathscr R_k\|y\|_{\rm original\ output}.
\tag{UPR26}
\]
This is a specified linear right inverse on the image.  Its values
may be passed to the quotient by \(\ker B_k\); it does not recover
components in that kernel.

\subsection{Original arithmetic quotient metric at every cutoff}
Now keep the original grid
\[
q=(k+1)^2,\quad c=k/2,\quad
\lambda_{a,b}=c+(2a-k)\delta+i(2b-k)\gamma,\quad
\chi(S)=\prod_{a,b=0}^k(S-\lambda_{a,b}),
\]
and let \(R=k\sqrt{\delta^2+\gamma^2}\),
\(d_*=2\min(\delta,\gamma)>0\).
Distinct roots have separation at least \(d_*\), because at least
one of their two grid-coordinate differences is a nonzero integer.
For an actual root-value column \(v\), the degree-\((q-1)\)
Lagrange representative is
\[
P_v(S)=\sum_\lambda
v_\lambda\frac{\chi(S)}{(S-\lambda)\chi'(\lambda)}.
\tag{UPR27}
\]
It has exactly those root values.  For \(S=c+iy\),
each numerator factor has modulus at most \(|y|+R\), whereas
\(|\chi'(\lambda)|\ge d_*^{q-1}\).
Cauchy--Schwarz in the actual \(q\)-term sum proves
\[
|P_v(c+iy)|^2
\le q\,d_*^{-2q+2}(|y|+R)^{2q-2}\|v\|_2^2.
\tag{UPR28}
\]

Retain the arithmetic source
\[
w_h(y)=|2\xi(1/2+iy)/h(1/2+iy)|^2/(2\pi),\qquad
m_k=w_h^{*k},
\quad
M_h(\eta)=\int_{\mathbb R}e^{\eta y}w_h(y)\,dy,\quad
\eta=\pi/4.
\]
The inherited exponential envelope gives a finite positive
\(M_h(\eta)\); no probability rescaling is made.
The nonnegative exponential series proves
\[
t^n\le n!\eta^{-n}e^{\eta t}\qquad(t\ge0,\ n\in\mathbb N).
\]
Tonelli's theorem in the original convolution coordinates gives
\(\int e^{\eta y}m_k(y)\,dy=M_h(\eta)^k\).
Evenness of the unchanged density yields
\[
\int_{\mathbb R}e^{\eta|y|}m_k(y)\,dy
=2\int_0^\infty e^{\eta y}m_k(y)\,dy
\le2M_h(\eta)^k.
\]
Apply these facts to (UPR28), including the factor \(e^{\eta R}\):
\[
\|P_v\|_{m_k}^2
\le\mathscr B_{h,k}\|v\|_2^2,\qquad
\mathscr B_{h,k}=
2q(2q-2)!(\eta d_*)^{-2q+2}
e^{\eta R}M_h(\eta)^k.
\tag{UPR29}
\]
Both half-lines, every root and every original source mass remain.

For each original cutoff \(N\ge q-1\), the minimum metric \(G_N\)
on the class of \(P_v\) is bounded above by the norm of this
particular admitted representative.  Composing (UPR26) with
(UPR27), and then with the actual class map, gives a linear
right inverse \(\mathscr T_{k,N}\) on the original observation image:
\[
\boxed{\displaystyle
\|\mathscr T_{k,N}\|_{{\rm original\ output}\to G_N}
\le \sqrt{\mathscr B_{h,k}}\,5^{-k/2}\mathscr R_k.}
\tag{UPR30}
\]
Passing its value further to the original quotient by
\(\ker\Lambda_k\) can only decrease the quotient norm, so the same
bound applies to that receiver.
Since \(\log(2q-2)!\le(2q-2)\log(2q-2)\), \(R=O_h(k)\),
and (UPR17) has logarithm \(O_{h,u}(k+\log(k+1))\),
\[
\log\|\mathscr T_{k,N}\|
\le O_{h,u}(q\log(q+2)),
\tag{UPR31}
\]
uniformly over every such \(N\).
All constants refer to the original fixed quartet, unit, period
and branch.  The construction is a large, finite recovery allowance
on the actual image.  It supplies no small native angular estimate,
does not remove the observation kernel, and does not identify the
arithmetic quotient metric with the independent proper-source metric.
\endgroup

\begin{thebibliography}{99}
\bibitem{human:stacks-trace}
The Stacks Project Authors, \emph{The Stacks Project},
Section 49.3, ``Discriminant of a finite locally free morphism'',
tag 0BVH, and Lemma 49.3.1, tag 0BJF.
\url{https://stacks.math.columbia.edu/tag/0BVH};
\url{https://stacks.math.columbia.edu/tag/0BJF}.
Read 19 September 2026. Used for the established multiplication-trace
construction; the exact original reflected tensor/residue factorization
is calculated in RBV31--35 from the original FTR/FRT providers.
\bibitem{programme:aks-current}
Split-Zero research programme,
\emph{Current cumulative mathematical source},
\texttt{09\_UPDATED\_JOINT\_NOTE.tex}, AKS7--10,
original lines 16755--16829.
The exact fixed-period matrix, root-value observation, ordinary
transpose convention and unscaled symmetric-word Gram are retained.

\bibitem{human:buchberger2006}
B. Buchberger,
\emph{Bruno Buchberger's PhD thesis 1965:
An algorithm for finding the basis elements of the residue class ring
of a zero dimensional polynomial ideal},
Journal of Symbolic Computation \textbf{41} (2006), 475--511,
English translation of the 1965 thesis.
\href{https://doi.org/10.1016/j.jsc.2005.09.007}
{doi:10.1016/j.jsc.2005.09.007}.
Bibliographic metadata checked against the RISC institutional
Gr\"obner Bases Bibliography, entry 712.

\bibitem{human:dickson1913}
L. E. Dickson,
\emph{Finiteness of the odd perfect and primitive abundant numbers
with \(n\) distinct prime factors},
American Journal of Mathematics \textbf{35} (1913), 413--422.
\href{https://doi.org/10.2307/2370405}{doi:10.2307/2370405}.
\bibitem{program:arrival03-01}
Supplied continuation, 19 September 2026,
\emph{The original coefficient image, its arithmetic compression,
and membership}, \texttt{01\_ORIGINAL\_COEFFICIENT\_IMAGE.md}.
Complete supplied bytes are preserved and pinned in
\texttt{SOURCE\_PINS.json}. COI1--17 and CJA give the exact
original-minor repair and subsequent full action calculation.

\bibitem{program:arrival03-08}
Supplied continuation, 19 September 2026,
\emph{All arithmetic coefficient-image flags: a uniform exponential
return estimate}, \texttt{08\_ALL\_COEFFICIENT\_ARITHMETIC\_FLAGS.md}.
CAL1--18 derives its whole-space estimate, strengthens its finite
tail bound, and records its exact image and original-minor scope.

\bibitem{program:arrival03-09}
Supplied continuation, 19 September 2026,
\emph{The full native-to-arithmetic coefficient metric return},
\texttt{09\_NATIVE\_ARITHMETIC\_METRIC\_RETURN.md}.
NMR1--25 derives the original relative determinant and stronger
uniform comparison cost while retaining its full arithmetic correction.

\bibitem{program:DNE}
Split-Zero research programme, \emph{Current kernel and multiplicity
proofs}, original DNE1--15 in
\texttt{14\_CURRENT\_KERNEL\_AND\_MULTIPLICITY\_PROOFS.tex},
source lines 20772--21116. DNE3 specifies the retained first
invertible minor in a fixed ordering of row subsets. The complete
source edition and its hash remain in \texttt{SOURCE\_PINS.json}.

\bibitem{program:NCT}
Split-Zero research programme, \emph{Current kernel and multiplicity
proofs}, NSO12 and NCT6--10, original source lines 49043--49073
and 63580--63708. The source fixes the full unscaled conductor,
the numerator moment map and the complete lower-evaluation minimum.
The upper-root lexicographic order and the row-subset enumeration
are retained as different specified pieces of data.

\bibitem{program:RTM}
Split-Zero research programme, \emph{The full tilted mass in the
original arithmetic tail} and \emph{A global tilted-mass transfer
with the central source retained}, complete supplied
\texttt{RTM\_GTM\_original\_providers.tex}, RTM1--11 and GTM1--21.
The original underlying convolution and tilt proofs are retained
as \texttt{THREE\_FACTOR\_ENVELOPE\_COMPLETE.tex} and
\texttt{BOUNDARY\_TILT\_COMPLETE.tex}.

\bibitem{program:GRA}
Split-Zero research programme, original GRA1--6, complete source
\texttt{14\_CURRENT\_KERNEL\_AND\_MULTIPLICITY\_PROOFS.tex},
lines 60782--60900. Its complete source constants and physical
coordinate are used literally in CAL1--4.

\bibitem{human:dlmf18}
T.~H. Koornwinder, R. Wong, R. Koekoek, R.~F. Swarttouw and
W.~P. Reinhardt, \emph{Orthogonal Polynomials}, Chapter 18 of the
\emph{NIST Digital Library of Mathematical Functions}.
General recurrence and zeros: \url{https://dlmf.nist.gov/18.2},
sections (iv) and (vi).
Classical Legendre orthogonality and Rodrigues formula:
\url{https://dlmf.nist.gov/18.3},
\url{https://dlmf.nist.gov/18.5.E5}.
Meixner--Pollaczek weight, norm and recurrence:
\url{https://dlmf.nist.gov/18.19},
\url{https://dlmf.nist.gov/18.22.E8}.
These primary references were consulted on 19 September 2026.
All finite recurrence and Legendre estimates needed here are
proved in the text with the original source mass.

\end{thebibliography}
